\documentclass[12pt]{article}

\usepackage{amsmath, amssymb, amsthm}
\usepackage{hyperref}
\usepackage{geometry}
\usepackage{booktabs}
\usepackage{graphicx}
\usepackage{draftwatermark}
\SetWatermarkText{WORKING PAPER}
\SetWatermarkScale{1.5}
\SetWatermarkAngle{45}
\SetWatermarkColor[gray]{0.80}
\geometry{margin=1in}

% Theorem environments
\newtheorem{theorem}{Theorem}[section]
\newtheorem{lemma}[theorem]{Lemma}
\newtheorem{proposition}[theorem]{Proposition}
\newtheorem{corollary}[theorem]{Corollary}
\newtheorem{conjecture}[theorem]{Conjecture}
\theoremstyle{definition}
\newtheorem{definition}[theorem]{Definition}
\theoremstyle{remark}
\newtheorem{remark}[theorem]{Remark}

% Notation
\newcommand{\vt}{v_2}
\newcommand{\Sm}{S}
\newcommand{\ZZp}{\mathbb{Z}^+}

\title{2-adic Valuations of $3n+1$ by Residue Class mod~8\\[0.5em]
       \large Speculative Results Series, Paper~66}

\author{Jon Seymour\\
  \texttt{jon@wildducktheories.com}\\
  Sydney, Australia}
\date{July 2026}

\begin{document}

\maketitle

\begin{abstract}
We prove that the 2-adic valuation $\vt(3n+1)$ under the Syracuse map
$\Sm(n) = (3n+1)/2^{\vt(3n+1)}$ is exactly determined — or has an exactly
determined expectation — by the residue class of $n$ modulo~8.
Specifically, writing $c_r = E[\vt(3n+1) \mid n \equiv r \pmod{8}]$:

\begin{center}
\begin{tabular}{ccc}
\toprule
$r$ & $c_r$ & Status \\
\midrule
1 & $2$ & exact constant \\
3 & $1$ & exact constant \\
5 & $4$ & exact expectation \\
7 & $1$ & exact constant \\
\bottomrule
\end{tabular}
\end{center}

For $r \in \{1,3,7\}$ the valuation is \emph{constant} on the entire residue
class — a direct consequence of the arithmetic of $3n+1$ modulo~8.
For $r = 5$ the valuation varies, but its expectation is exactly~4:
writing $n = 8k+5$ gives $\vt(3n+1) = 3 + \vt(3k+2)$, and
$E[\vt(3k+2)] = 1$ follows from $\gcd(3, 2^j) = 1$, which makes $3k+2$
uniformly distributed modulo~$2^j$ for every $j \geq 1$.

The sum of the four expected valuations is $c_1 + c_3 + c_5 + c_7 = 8 = 4 \times 2$,
so the mean valuation across all four classes equals~$2$, the same as the
constant valuation of class~1. Equivalently, $3^4/2^8 = 81/256 = (3/4)^4$:
the product $\prod_r 3/2^{c_r}$ of the four asymptotic contraction ratios
equals $(3/4)^4$, whose fourth root is exactly $3/4$.

This result provides the exact integer-exponent foundation for understanding
the asymptotic contraction rates of the Syracuse map by residue class, and
in particular why the $5\pmod{8}$ class — with $c_5 = 4$, giving asymptotic
ratio $3/2^4 = 3/16$ — acts as a dynamical sink that drives orbits toward~1
with exceptional efficiency.

\medskip
\noindent\textbf{Status.} This paper contains proved theorems. The main results
(Theorem~\ref{thm:v2} and Corollary~\ref{cor:balance}) are unconditional.
The dynamical interpretation (Section~\ref{sec:dynamics}) is stated as
observation and conjecture.
\end{abstract}

\tableofcontents

% ---------------------------------------------------------------------------
\section{Background and Notation}
\label{sec:background}

We work with odd positive integers throughout. The \emph{2-adic valuation}
$\vt(n)$ denotes the exponent of the largest power of~$2$ dividing~$n$.

\begin{definition}[Syracuse map]
\label{def:syracuse}
The \emph{Syracuse map} $\Sm : \mathbb{O} \to \mathbb{O}$ on odd positive
integers is
\[
  \Sm(n) \;=\; \frac{3n+1}{2^{\vt(3n+1)}}.
\]
\end{definition}

The ratio $\Sm(n)/n = (3n+1)/(n \cdot 2^{\vt(3n+1)})$ measures the multiplicative
change per Syracuse step. Its asymptotic behaviour as $n \to \infty$ is governed
entirely by the exponent $\vt(3n+1)$, since $(3n+1)/n \to 3$. The central
question of this paper is: what is $\vt(3n+1)$, and how does it depend on
$n \bmod 8$?

% ---------------------------------------------------------------------------
\section{Main Theorem: Valuations by Residue Class}
\label{sec:main}

\begin{theorem}[2-adic valuations of $3n+1$ mod~8]
\label{thm:v2}
Let $n$ be an odd positive integer and write $c_r = E[\vt(3n+1) \mid n \equiv r \pmod{8}]$
where the expectation is over the uniform (natural) density on the residue class.
Then:
\begin{enumerate}
  \item $n \equiv 1 \pmod{8}$: $\vt(3n+1) = 2$ for every such $n$, so $c_1 = 2$.
  \item $n \equiv 3 \pmod{8}$: $\vt(3n+1) = 1$ for every such $n$, so $c_3 = 1$.
  \item $n \equiv 5 \pmod{8}$: $\vt(3n+1) \geq 3$ for every such $n$,
        and $E[\vt(3n+1) \mid n \equiv 5 \pmod{8}] = 4$, so $c_5 = 4$.
  \item $n \equiv 7 \pmod{8}$: $\vt(3n+1) = 1$ for every such $n$, so $c_7 = 1$.
\end{enumerate}
\end{theorem}

\begin{proof}
Write $n = 8k + r$ for $r \in \{1,3,5,7\}$ and $k \geq 0$.

\medskip\noindent
\textbf{(1) $r = 1$:} $3n+1 = 3(8k+1)+1 = 24k+4 = 4(6k+1)$.
Since $6k$ is even, $6k+1$ is odd, so $\vt(6k+1) = 0$ and
$\vt(3n+1) = \vt(4) + \vt(6k+1) = 2 + 0 = 2$ for every $k$.

\medskip\noindent
\textbf{(2) $r = 3$:} $3n+1 = 3(8k+3)+1 = 24k+10 = 2(12k+5)$.
Since $12k$ is even, $12k+5$ is odd, so $\vt(3n+1) = 1$ for every $k$.

\medskip\noindent
\textbf{(3) $r = 5$:} $3n+1 = 3(8k+5)+1 = 24k+16 = 8(3k+2)$.
Hence $\vt(3n+1) = 3 + \vt(3k+2)$ for every $k$, establishing $\vt(3n+1) \geq 3$.
It remains to show $E[\vt(3k+2)] = 1$.

For any $j \geq 1$, since $\gcd(3, 2^j) = 1$, multiplication by~3 is a
bijection on $\mathbb{Z}/2^j\mathbb{Z}$. Therefore as $k$ ranges uniformly
over $\mathbb{Z}/2^j\mathbb{Z}$, the value $3k+2 \pmod{2^j}$ is also uniform
over $\mathbb{Z}/2^j\mathbb{Z}$. In particular:
\[
  P\!\left(2^j \mid 3k+2\right) \;=\; \frac{1}{2^j}.
\]
Using the standard identity for non-negative integer-valued random variables:
\[
  E[\vt(3k+2)] \;=\; \sum_{j=1}^{\infty} P\!\left(\vt(3k+2) \geq j\right)
  \;=\; \sum_{j=1}^{\infty} P\!\left(2^j \mid 3k+2\right)
  \;=\; \sum_{j=1}^{\infty} \frac{1}{2^j} \;=\; 1.
\]
Therefore $c_5 = E[\vt(3n+1) \mid n \equiv 5 \pmod{8}] = 3 + 1 = 4$.

\medskip\noindent
\textbf{(4) $r = 7$:} $3n+1 = 3(8k+7)+1 = 24k+22 = 2(12k+11)$.
Since $12k+11$ is odd ($12k$ even, $+11$ odd), $\vt(3n+1) = 1$ for every $k$.
\end{proof}

\begin{remark}
Parts (1), (2), (4) are exact pointwise statements: the valuation is a fixed
integer on the entire residue class. Part (3) is an expectation statement: the
valuation varies over the class $n \equiv 5 \pmod{8}$, but its mean is exactly~4.
The key tool in part~(3) is the invertibility of multiplication by~3 modulo~$2^j$,
i.e.\ $\gcd(3, 2^j) = 1$, which makes the distribution of $3k+2 \pmod{2^j}$
exactly uniform.
\end{remark}

\begin{remark}
Theorem~\ref{thm:v2} explains the mod-8 step taxonomy of~\cite{paper33}:
classes 1, 3, 7 each have a fixed valuation (step types OEE, OE, OE respectively),
while class~5 has a variable valuation $\geq 3$ (step type OEEE$^+$, at least
three even steps).
\end{remark}

% ---------------------------------------------------------------------------
\section{The Balance Identity}
\label{sec:balance}

\begin{corollary}[Balance of expected valuations]
\label{cor:balance}
The expected valuations satisfy:
\[
  c_1 + c_3 + c_5 + c_7 \;=\; 2 + 1 + 4 + 1 \;=\; 8 \;=\; 4 \times 2.
\]
The mean expected valuation across all four classes is exactly~$2 = c_1$,
the constant valuation of the neutral class~$1 \pmod{8}$.

Equivalently, in terms of powers of~$2$ and~$3$:
\[
  \frac{3^4}{2^{c_1} \cdot 2^{c_3} \cdot 2^{c_5} \cdot 2^{c_7}}
  \;=\; \frac{3^4}{2^8}
  \;=\; \frac{81}{256}
  \;=\; \left(\frac{3}{4}\right)^4.
\]
\end{corollary}

\begin{proof}
Direct: $2 + 1 + 4 + 1 = 8$ and $2^{c_1}\cdot 2^{c_3}\cdot 2^{c_5}\cdot 2^{c_7}
= 2^{2+1+4+1} = 2^8 = 256$.
\end{proof}

\begin{remark}
The balance identity has a clean interpretation in terms of the asymptotic
contraction ratios. For large $n$ in class $r$, the Syracuse step multiplies $n$
by approximately $3/2^{c_r}$:
\[
  \frac{\Sm(n)}{n} \;=\; \frac{3n+1}{n \cdot 2^{\vt(3n+1)}}
  \;\longrightarrow\; \frac{3}{2^{c_r}} \quad \text{as } n \to \infty.
\]
The four asymptotic ratios are $3/4$, $3/2$, $3/16$, $3/2$ for classes
$1, 3, 5, 7$ respectively. Their product is $3^4/2^8 = (3/4)^4$, whose fourth
root is $3/4$. That is: the \emph{geometric mean} of the four asymptotic
contraction ratios equals the ratio of the neutral class~1.

Note that the ratios $3/2^{c_r}$ are exact as limits ($n \to \infty$) for
classes 1, 3, 7 where $\vt$ is constant, and exact as expectations for class~5.
The finite-$n$ correction is $O(1/n)$ in each case and does not affect the
integer-exponent result.
\end{remark}

\begin{remark}[Asymmetry of the four classes]
The four classes play structurally distinct roles:
\begin{itemize}
  \item Classes 3 and 7 are \emph{growth} classes: $c_3 = c_7 = 1$,
        asymptotic ratio $3/2 > 1$.
  \item Class 1 is the \emph{neutral} class: $c_1 = 2$, asymptotic ratio $3/4 < 1$.
  \item Class 5 is the \emph{contraction} class: $c_5 = 4$, asymptotic ratio
        $3/16 \ll 1$. A single step in class~5 contributes $c_5 = 4$ to the
        sum of valuations, the same as \emph{four} steps in classes 3 or 7,
        or \emph{two} steps in class~1.
\end{itemize}
The balance $c_1 + c_3 + c_5 + c_7 = 4 c_1$ is equivalent to saying
class~5's excess valuation $(c_5 - c_1) = 2$ exactly compensates the
deficit of classes 3 and 7 combined: $(c_1 - c_3) + (c_1 - c_7) = 2$.
\end{remark}

% ---------------------------------------------------------------------------
\section{Dynamical Implications}
\label{sec:dynamics}

\subsection{Class 5 as a Dynamical Sink}

Theorem~\ref{thm:v2} gives a precise, quantitative sense in which $5\pmod{8}$
acts as a \emph{dynamical sink} in the Collatz tree. Each visit to a $5\pmod{8}$
node contributes an expected valuation of~4 to the denominator of the Syracuse
ratio, compared to~2 for class~1, and only~1 for classes 3 and 7.
In the asymptotic ratio sense, class~5 contracts by $3/2^4 = 3/16$ per step
— four times the exponent of class~1 ($3/2^2 = 3/4$) and sixteen times the
inverse of classes 3 and 7 ($3/2^1 = 3/2$).

\subsection{Connection to 5-Dominance in Collatz Paths}

It is empirically observed that of the first million Collatz sequences, 93.8\%
pass through the integer~5 on the way to~1, with only 6.2\% passing through~32
instead~\cite{reddit-observation}. Theorem~\ref{thm:v2} gives the structural
foundation: any orbit entering the $5\pmod{8}$ class experiences a valuation
of~4 in expectation, driving it toward~1 far more rapidly than any other class.
The complementary probability 6.2\% $\approx 1/16$ corresponds to the mod-32
edge density of $5\to5$ transitions in the Syracuse graph (see Paper~65
\cite{paper65}).

\subsection{Connection to Steiner Sentence Length Distribution}

Paper~65~\cite{paper65} empirically observes that Steiner sentence lengths
follow $P(k) = \frac{1}{3}(3/4)^k$ rather than the naive $(1/2)^k$.
The balance identity $c_1 + c_3 + c_5 + c_7 = 4 c_1$ is the exact integer
expression of the $3/4$ decay rate: the neutral class~1 (which governs
within-sentence dynamics after conditioning on non-termination) has valuation
exactly $c_1 = 2$, giving asymptotic ratio $3/4$. The first-principles proof
is given in Paper~67~\cite{paper67}.

A key insight of Paper~67 is that Steiner sentences are the \emph{right scale}
for this analysis precisely because they make class~5 a sink rather than a source.
Within a sentence, the chain visits only classes~1, 3, and~7 until it terminates
at class~5. For these three classes, the 2-adic valuation $\vt(3n+1)$ is a
\emph{constant} (2, 1, 1 respectively — Theorem~\ref{thm:v2}), so the transition
probabilities are exact rationals with no averaging required.
The variable-valuation behaviour of class~5 — the central subject of this paper,
requiring the geometric series argument of Theorem~\ref{thm:v2}(3) — never appears
in the interior of a sentence. It enters only at the sentence boundary (the
uniform output of class~5 that starts each new sentence) and as the termination
condition (class~5 is the absorbing state). In both roles, only row~5 of the
transition matrix matters, and that row is uniform by the same $\gcd(3,2^j)=1$
argument.
The two papers are thus complementary: this paper characterises the hard case
(class~5, unbounded valuation); Paper~67 sidesteps it by choosing the right scale.

% ---------------------------------------------------------------------------
\section{Further Results}
\label{sec:open}

\begin{proposition}[Full distribution for class~5]
\label{prop:class5dist}
For $n \equiv 5 \pmod{8}$ drawn with natural density,
\[
  P\!\bigl(\vt(3n+1) = j \;\big|\; n \equiv 5 \pmod{8}\bigr)
  \;=\; \frac{1}{2^{j-2}}, \qquad j \geq 3.
\]
\end{proposition}

\begin{proof}
Writing $n = 8k+5$ gives $\vt(3n+1) = 3 + \vt(3k+2)$
(as in Theorem~\ref{thm:v2}(3)).
The pmf of $\vt(3k+2)$ under natural density is
$P(\vt(3k+2) = m) = 1/2^{m+1}$ for $m \geq 0$:
this follows from $P(\vt(3k+2) \geq m) = P(2^m \mid 3k+2) = 1/2^m$
(by the uniformity argument of Theorem~\ref{thm:v2}(3)) via
$P(\vt = m) = P(\vt \geq m) - P(\vt \geq m+1) = 1/2^m - 1/2^{m+1} = 1/2^{m+1}$.
Setting $j = 3 + m$ gives $P(\vt(3n+1) = j) = 1/2^{(j-3)+1} = 1/2^{j-2}$.
Normalisation: $\sum_{j \geq 3} 1/2^{j-2} = \sum_{i \geq 1} 1/2^i = 1$. \qed
\end{proof}

\begin{remark}
This proposition was stated as Open Problem~2 in a draft of this paper.
It was machine-verified in Lean~4/Mathlib by the Aristotle automated
proof assistant (July 2026); see \texttt{ARISTOTLE-REVIEW-66.md} and
\texttt{Paper66.lean} in the paper directory.
\end{remark}

\begin{enumerate}
  \item \textbf{Higher moduli.} The argument for class~5 uses $\gcd(3,2^j)=1$
        to establish uniform distribution of $3k+2 \pmod{2^j}$. Do analogous
        exact valuation results hold for finer residue classes modulo~16, 32,
        and beyond? The mod-32 structure of $5\to5$ edges in the Syracuse
        graph suggests a rich mod-32 refinement.

  \item \textbf{First-principles sentence length distribution.}
        The bridge from Theorem~\ref{thm:v2} to a proof of
        $P(\text{sentence length} = k) = \frac{1}{3}(3/4)^k$ is given
        in Paper~67~\cite{paper67}.
\end{enumerate}

% ---------------------------------------------------------------------------
\section*{Acknowledgements}

This paper is part of the speculative results stream (64+k) of the Collatz
programme. Paper~33~\cite{paper33} provides the Steiner circuit framework.
Paper~65~\cite{paper65} provides the empirical motivation.

% ---------------------------------------------------------------------------
\begin{thebibliography}{9}

\bibitem{paper33}
J.~Seymour,
\emph{A Regular Expression Language for the Collatz Graph} (Working Paper, Paper~33),
2026.

\bibitem{paper65}
J.~Seymour,
\emph{The $3^{k-1}/4^k$ Distribution of Steiner Sentence Lengths}
(Working Paper, Paper~65),
2026.

\bibitem{paper67}
J.~Seymour,
\emph{First-Principles Derivation of the Steiner Sentence Length Distribution}
(Working Paper, Paper~67),
2026.

\bibitem{reddit-observation}
Reddit post (r/Collatz),
\emph{Why does 5 dominate Collatz paths to 1?}, 2026.

\end{thebibliography}

\end{document}
